\section{Introduction}



% We will now discuss 12d perspectives on the type IIB effective action, and argue that type IIB effective actions can not be written in 12d covariant form.
% This is separate from the arguments related to the Eisenstein series structure, as well as the U(1)-violating sector, which has been recognized as difficult to be understood from 12d perspective \cite{Minasian:2015bxa,Liu:2025uqu}.
% However, nevertheless, one may still hope that 12d contractions, upon reduction to 10d, and after applying LSZ reduction, produce the kinematic structure of string scattering amplitudes.
% We argue that this can not be done.



% A general strategy for uplifting type IIB effective action may look like the following. Examine the effective action, substituting $\tau_1, \tau_2$ for components of the 12d metric on $T_2$, and derivatives of $\tau_1, \tau_2$ for 12d Riemannt tensor.
% This packages SL(2, Z) as part of 12d covariance. 
% One may then examine whether the resulting expression, in particular the coefficient of the distinct contraction patterns, may have a natural interpretation as arising from a higher dimensional contraction.
% The idea is that if certain terms of the 10d effective action arises from a 12d contraction, then one must be able to rewrite the 10d effective action to invert the index splitting, and promote a group of 10d terms to a 12d contraction.


% We note that crucially, while for a generic 12d metric ansatze, $R_{ambn}, R_{mnab}, R_{abcd}$ are independent of each other, for the specific F-theory ansatze, $R_{mnab}$ and $R_{abcd}$ can both be obtained as linear combinations or contractions of $R_{ambn}$.
% For example, $R_{z\overline{z}z\overline{z}} = \tfrac12 R^m{}_{zm\overline{z}}$, and $R_{mnz\overline{z}} = R_{m\overline{z}nz} - R_{mzn\overline{z}}$, leaving us with only 4 linearly independent components.
% This basis mismatch is related to the specific 12d metric ansatze chosen, and has important consequences as we will see.

